\section{Global uniqueness on a roof--angle fiber}
\label{sec:v111-global}
The estimates in this section concern actual finite trajectories,
not a continuation through a forbidden event. In particular the
regular domain of a lifted itinerary need not be connected or convex.
This distinction is important when the original section and occupation
tests divide that domain into many components.

Normalize the initial lifted scatterer to $O_0=0$. Let
$\mathcal W_m(R)$ denote the realizable lifted itineraries of $m$
flights. This set is finite by the finite-horizon bound and the lattice
geometry. Fix $w=(O_0,\ldots,O_m)\in\mathcal W_m(R)$. For a physical trajectory
write $Q_j=O_j+Rn_j$, let $t_j$ be the counterclockwise rotation of
$n_j$, and put
\[
 S(Q)=\sum_{j=0}^{m-1}|Q_{j+1}-Q_j|,\qquad
 g=R(-p_0,p_m),\qquad z=|g|.
\]
As a function of the independent Cartesian contacts, $S$ is convex.
It is differentiable at every configuration in question because
successive scatterer disks are disjoint. Write $\beta_0=\beta_m=1$
and $\beta_j=2$ for $0<j<m$.

\begin{lemma}[Two-trajectory support inequalities]
\label{lem:v111-support}
Let $Q$ and $\widehat Q$ be two regular physical trajectories with
the same lifted itinerary and radius. Set
\[
 b_j=1-n_j\cdot\widehat n_j\ge0,\qquad
 s_j=t_j\cdot\widehat n_j,
 \qquad s_E=(s_0,s_m).
\]
Their lengths satisfy
\begin{align}
 S(\widehat Q)-S(Q)
 &\ge R\sum_{j=0}^m\beta_jc_jb_j+g\cdot s_E,
                                      \label{eq:v111-support-forward}\\
 S(Q)-S(\widehat Q)
 &\ge R\sum_{j=0}^m\beta_j\widehat c_jb_j
                                      -\widehat g\cdot s_E.
                                      \label{eq:v111-support-backward}
\end{align}
Consequently
\begin{equation}\label{eq:v111-support-sum}
 (\widehat g-g)\cdot s_E
 \ge R\sum_{j=0}^m\beta_j(c_j+\widehat c_j)b_j.
\end{equation}
No interpolation by physical trajectories is used.
\end{lemma}
\begin{proof}
The physical reflection law and first variation in Cartesian
coordinates give
\[
 \nabla_0S=-c_0n_0-p_0t_0,\qquad
 \nabla_jS=-2c_jn_j\ (0<j<m),\qquad
 \nabla_mS=-c_mn_m+p_mt_m.
\]
The convex support inequality is
$S(\widehat Q)-S(Q)\ge\sum_j\nabla_jS(Q)\cdot(\widehat Q_j-Q_j)$.
Since $\widehat Q_j-Q_j=R(\widehat n_j-n_j)$, its normal terms
are $R\beta_jc_jb_j$ and its two tangential terms are $g\cdot s_E$.
This proves \eqref{eq:v111-support-forward}. Reversing the roles of
the two configurations leaves $b_j$ unchanged and replaces $s_j$
by $-s_j$. This proves \eqref{eq:v111-support-backward}.
Adding the two inequalities proves \eqref{eq:v111-support-sum}.
The support inequality is applied to a convex Cartesian function;
it requires neither a convex angular chart nor a physical path
joining the configurations.
\end{proof}

\begin{theorem}[Global roof--angle uniqueness]
\label{thm:v111-unique}
For a fixed lifted itinerary, radius and direction $e\in S^1$,
the physical length is strictly increasing among regular trajectories
whose endpoint gradients are $g=ze$, $z>0$, ordered by $z$.
At a fixed $z$ there is at most one such trajectory. In particular
\begin{equation}\label{eq:v111-map}
                    Q\longmapsto(F(Q),\arg g(Q))
\end{equation}
is globally injective on the noncritical regular trajectories of
that itinerary. This remains true after any of the original exact
event tests, and across all their disconnected regular components.
\end{theorem}
\begin{proof}
If $\widehat g=g$, the left side of
\eqref{eq:v111-support-sum} is zero. All $c_j$ and $\widehat c_j$
are strictly positive, so every $b_j=0$. Thus all physical contacts,
and then all flight directions, coincide.

Now let $\widehat g=\widehat z e$, $g=ze$ and $\widehat z>z>0$.
The configurations cannot coincide. At least one $b_j$ is therefore
positive, and \eqref{eq:v111-support-sum} gives
\[
 (\widehat z-z)e\cdot s_E
       \ge R\sum_j\beta_j(c_j+\widehat c_j)b_j>0.
\]
Hence $e\cdot s_E>0$. The right side of
\eqref{eq:v111-support-forward} is then strictly positive, proving
$S(\widehat Q)>S(Q)$. Equal roof and angle would force equal $z$,
and the first part of the proof would force equality of trajectories.
Restricting an injective map to the original exact event cannot
increase its multiplicity.
\end{proof}

\begin{lemma}[A bounded roof range for each physical itinerary]
\label{lem:v111-roof-range}
For any two regular physical trajectories with the same itinerary,
\begin{equation}\label{eq:v111-roof-range}
 |S(Q)-S(\widehat Q)|
       \le |Q_0-\widehat Q_0|+|Q_m-\widehat Q_m|\le4R.
\end{equation}
\end{lemma}
\begin{proof}
With its endpoint contacts fixed, a regular physical trajectory
minimizes $S$ over the product of the closed interior scatterer
disks. Indeed, for a competitor $Y_j$ in the $j$th disk,
$n_j\cdot(Y_j-Q_j)\le0$; the convex support inequality and
$\nabla_jS=-2c_jn_j$ therefore give the minimizing property.
Use the interior contacts of $\widehat Q$ as a competitor for the
endpoints of $Q$. Changing only the first and last endpoint changes
the sum of flight lengths by at most the sum of their displacements,
by the triangle inequality. This proves one direction of
\eqref{eq:v111-roof-range}; interchanging the trajectories proves
the other. The same triangle inequality covers $m=1$.
\end{proof}
This range estimate does not control how many original occupation
labels occur in a word domain. In particular it is not used to
replace the exact occupation test by an unqualified interval test.

\subsection{The original physical sum at a target roof}
For the original section tests, $\chi\ge\lambda_\circ$ by
\eqref{eq:v109-chi}. On each regular component, the map
\eqref{eq:v111-map} is a local diffeomorphism. Theorem~\ref{thm:v111-unique}
shows that all those components of a fixed itinerary have disjoint
images, not merely that each component has a nonsingular Jacobian.
The equal-gradient argument also gives at most one regular critical
trajectory per word, even when the gradient is zero. There are
finitely many physical words at each finite count, so these critical
points have zero collision area. Consequently, for any nonnegative
Borel trajectory weight $a$, the
original roof density has the representation
\begin{equation}\label{eq:v111-roof-sum}
 \frac{dF_\#(a\one_\Pi\nu_R^*)}{dt}(t)
  =\int_0^{2\pi}\sum_{w\in\mathcal W_m(R)}
       \one_{\{x_w(t,\vartheta)\ \mathrm{exists}\}}
       \frac{\omega(x_w)a(x_w)}{4\pi c_*\chi(x_w)}\,d\vartheta.
\end{equation}
Here $x_w$ is the unique point of the original exact event with the
indicated word, roof and angle. A missing point contributes zero.
The representation holds for almost every roof, in the original
order $\sup_R\esssup_t$. Angular chart overlap is removed before
summation; the incident-component convention remains unchanged.

\begin{proposition}[Monotone completion of the target-roof formula]
\label{prop:v112-exhaustion}
Formula~\eqref{eq:v111-roof-sum} is an identity of nonnegative
measures for every nonnegative Borel weight, with extended values
allowed before an integrability bound is proved. In particular it
applies to the unbounded products $\prod_{i\in S}c_i^{-1}$.
\end{proposition}
\begin{proof}
Fix $m,R,n,k$. Cover the noncritical regular physical domain in
endpoint coordinates by countably many relatively compact local
charts $V_j$ on which $(F,\arg g)$ is a $C^1$ diffeomorphism into
an unwrapped angular chart. Such a cover exists by the nonvanishing
Jacobian and second countability. Subdivide as necessary so that
both maps have bounded derivatives on each smaller chart.
Disjointize the source charts by
$B_j=V_j\setminus\bigcup_{k<j}V_k$. The sets $B_j$ are Borel;
no chart overlap is counted twice. Angular representatives are also
split at a fixed seam, so their range is $[0,2\pi)$.

For a nonnegative Borel weight $h$ put
\[
 h_N=\min\{h,N\}\one_{\bigcup_{j\le N}B_j}.
\]
This sequence increases to $h$ on the regular noncritical domain.
Apply ordinary change of variables on each of the first $N$
charts, restricted to its Borel piece, and test against a
nonnegative Borel function of $(t,\vartheta)$. The local inverse
Jacobian, multiplied by the original collision density, is exactly
$\omega/(4\pi c_*\chi)$. Local inverses are continuous and agree
where they represent the same physical point, so their Borel
restrictions define the displayed $x_w$ without a selection ambiguity.
Global uniqueness prevents different regular components of one word
from producing the same target. Distinct words are summed, not
identified.

Monotone convergence on the source, followed by Tonelli and monotone
convergence on the target, now gives the identity for $h$ as an
extended nonnegative measure. Only after this passage is an upper
pressure bound used to establish finiteness and an $L^\infty$
density. This avoids assuming the integrability of the unbounded
marked weight in order to prove its own change of variables.

The complement of the two-sided regular set has zero collision
area. On the countably many charts just chosen, absolute continuity
of the coordinate changes and their local Lipschitz property give
zero target area for its image. The finitely many regular critical
points also have zero source area. Removing these sets therefore
does not change the weighted source measure. Fubini yields the
claimed roof identity for almost every $t$. All null sets here are
chosen at the fixed radius; no exceptional set common to every
$R$ is asserted or needed. For the finitely many marked sets at a
fixed $m$, intersecting the corresponding full-measure roof sets
provides a common version for their finite positive expansions.
\end{proof}
